\documentclass[11pt]{amsart}
\usepackage[margin=1.05in]{geometry}
\usepackage{amsmath,amssymb,amsthm,mathtools}
\usepackage{booktabs}
\usepackage[shortlabels]{enumitem}
\usepackage[expansion=false]{microtype}
\usepackage{xcolor}
\usepackage[colorlinks=true,linkcolor=blue!50!black,citecolor=blue!50!black,urlcolor=blue!50!black]{hyperref}

\hypersetup{pdftitle={Proof of a conjecture of Saikia and Talukdar on internal 2-adic congruences for overpartition tuples},pdfauthor={Samson Fraser}}

\newtheorem{theorem}{Theorem}[section]
\newtheorem{proposition}[theorem]{Proposition}
\newtheorem{lemma}[theorem]{Lemma}
\newtheorem{corollary}[theorem]{Corollary}
\newtheorem{conjecture}[theorem]{Conjecture}
\theoremstyle{definition}
\newtheorem{definition}[theorem]{Definition}
\theoremstyle{remark}
\newtheorem{remark}[theorem]{Remark}

\newcommand{\Z}{\mathbb{Z}}
\newcommand{\Q}{\mathbb{Q}}
\newcommand{\OPT}{\overline{\mathrm{OPT}}}
\newcommand{\Sop}{\mathcal{S}}
\newcommand{\Th}{\widetilde{\Theta}}
\newcommand{\Dc}{\mathcal{D}}

\title[Proof of a conjecture of Saikia and Talukdar]{Proof of a conjecture of Saikia and Talukdar on internal $2$-adic congruences for overpartition tuples}
\author{Samson Fraser}
\address{Independent researcher, Sicamous, British Columbia, Canada}
\email{samsonmfraser@gmail.com}
\subjclass[2020]{Primary 11P83; Secondary 05A17}
\keywords{overpartitions, overpartition tuples, internal congruences, $2$-adic valuations, Chebyshev polynomials}
\date{}



\begin{document}
\begin{abstract}
Let $\OPT_m(n)$ be the number of overpartition $m$-tuples of $n$ into odd parts. For even $m$, Saikia and Talukdar reduced the internal congruences $\OPT_m(2^in)\equiv\OPT_m(2^{i-1}n)$ to a conjectural pattern in the $2$-adic valuations of a family of integer polynomials $R^{(m)}_i$ (their Conjecture~6.2), and they deduced congruences modulo $2^{3i+c_m}$ from it (their Corollary~6.3). We prove the conjecture for every even $m\ge 2$ and every $i>i_m$. We also show that, for $i>i_m$, $\nu(R^{(m)}_i)=3i-8$ if $m\equiv 2\pmod 4$ and $\nu(R^{(m)}_i)=3i+2\nu(m)-9$ if $4\mid m$. Consequently $c_m=-4$ when $\nu(m)=1$ and $c_m=2\nu(m)-5$ when $\nu(m)\ge 2$. This confirms the authors' suggestion that $c_m$ depends only on $\nu(m)$, and it makes their Corollary~6.3 unconditional.
\end{abstract}
\maketitle

\section{Introduction}

For $m\ge1$ let $\OPT_m(n)$ denote the number of overpartition $m$-tuples of $n$ whose parts are all odd. Its generating function is
\[
  \sum_{n\ge 0}\OPT_m(n)q^n \;=\; A(q)^m,\qquad
  A(q)=\prod_{j\ge1,\ j\ \mathrm{odd}}\frac{1+q^j}{1-q^j}=\frac{f_2^3}{f_1^2f_4},\qquad f_k=(q^k;q^k)_\infty .
\]
Saikia and Talukdar~\cite{ST} proved, for every odd $m$ and every $i\ge3$, the congruence
\[
  \OPT_m(2^in)\equiv\OPT_m(2^{i-1}n)\pmod{2^{i+1}},
\]
and that the $2$-adic valuation of the difference equals $i+1$ exactly at the odd squares. In~\cite[\S6]{ST} they treat even $m$ through the following reduction, which we recall in their notation except that we write $T_n$ for the Chebyshev polynomials of the first kind (they write $D_n$):
\[
  T_0=1,\qquad T_1=x,\qquad T_{n+1}=2xT_n-T_{n-1}.
\]

\subsection*{The setting of \texorpdfstring{\cite[\S6]{ST}}{[ST, Sec. 6]}}
Let $U\big(\sum c(n)q^n\big)=\sum c(2n)q^n$ and $\xi=A^4$. Then $\xi\in 1+8q\Z[[q]]$ by \cite[(2.16)]{ST} (see also Lemma~\ref{lem:Uiden}), and in fact $\xi=1+8q+O(q^2)$ because $A=1+2q+O(q^2)$. Let $\Sop\colon\Z[x]\to\Z[x]$ be the $\Z$-linear map with $\Sop(x^j)=T_{2j}(x)$, and put
\[
  i_m=\begin{cases}1,&4\mid m,\\ 2,&m\equiv2\pmod 4,\end{cases}\qquad
  P_{i_m-1}=\begin{cases}x^{m/4},&4\mid m,\\ T_{m/2}(x),&m\equiv2\pmod 4,\end{cases}\qquad
  P_{i+1}=\Sop(P_i).
\]
Saikia and Talukdar show that $U^i(A^m)=P_i(\xi)$ for $i\ge i_m-1$ \cite[\S6, (6.3)--(6.4) and the display following (6.4)]{ST}; for completeness we include a short proof in Lemma~\ref{lem:Uiden}. Consequently, with
\[
  R_i=R^{(m)}_i:=\frac{P_i-P_{i-1}}{x-1}\in\Z[x]\qquad(i\ge i_m),
\]
the generating function of the differences is
\begin{equation}\label{eq:65}
  \Dc^{(m)}_i:=\sum_{n\ge0}\big(\OPT_m(2^in)-\OPT_m(2^{i-1}n)\big)q^n=(\xi-1)\,R^{(m)}_i(\xi)\qquad(i\ge i_m),
\end{equation}
which is \cite[(6.5)]{ST}. For a nonzero integer $c$ let $\nu(c)$ be its $2$-adic valuation, and put $\nu(0)=\infty$. For $F\in\Z[x]$ let $\nu(F)$ be the minimum of $\nu$ over the coefficients of $F$ in the monomial basis. Thus $\nu(F)\ge r$ if and only if $F\in 2^r\Z[x]$.

\cite[Lemma~6.1]{ST} introduces the operator
\[
  \Th(R):=\frac{\Sop\big((x-1)R\big)}{x-1}\qquad(R\in\Z[x]),
\]
shows that $\Th(R)\in\Z[x]$ and $\Th(R)(1)=4R(1)+8R'(1)$, and shows that $R_{i+1}=\Th(R_i)$ for $i\ge i_m$.

\begin{conjecture}[{\cite[Conjecture~6.2]{ST}}]\label{conj}
Let $m\ge2$ be even, and put $v_i:=\nu(R^{(m)}_i)$. For every $i>i_m$:
\[
  \nu\big((R^{(m)}_i)'(1)\big)=v_i,\qquad \nu\big(R^{(m)}_i(1)\big)=v_i+1,\qquad v_{i+1}=v_i+3 .
\]
\end{conjecture}

Assuming Conjecture~\ref{conj}, \cite[Corollary~6.3]{ST} states that for each even $m$ there is an integer $c_m$ with
\begin{equation}\label{eq:cong}
  \OPT_m(2^in)\equiv\OPT_m(2^{i-1}n)\pmod{2^{3i+c_m}}\qquad(i>i_m,\ n\ge0),
\end{equation}
where $3i+c_m=v_i+4$ and the modulus is attained at $n=1$. Computations in \cite{ST}, for $2\le m\le14$ and $i\le 9$ and carried out modulo $2^{45}$, support the conjecture. They give $c_m=-4$ for $m\in\{2,6,10,14\}$, $c_4=c_{12}=-1$ and $c_8=1$, and the authors suggest that $c_m$ may depend only on $\nu(m)$.

\subsection*{Results}

\begin{theorem}\label{thm:main}
Conjecture~\ref{conj} holds for every even $m\ge2$ and every $i>i_m$. Moreover, for $i>i_m$,
\[
  v_i=\nu\big(R^{(m)}_i\big)=\begin{cases}3i-8,&m\equiv 2\pmod 4,\\[2pt] 3i+2\nu(m)-9,&4\mid m.\end{cases}
\]
\end{theorem}

\begin{corollary}\label{cor:cm}
For every even $m\ge2$, the constant of \cite[Corollary~6.3]{ST} is
\[
  c_m=\begin{cases}-4,&\nu(m)=1,\\ 2\nu(m)-5,&\nu(m)\ge2.\end{cases}
\]
In particular $c_m$ depends only on $\nu(m)$. Since the identity~\eqref{eq:65} does not depend on the conjecture (it is proved in \cite[\S6]{ST}, and again in Lemma~\ref{lem:Uiden} below), the argument recalled in Section~\ref{sec:proof} shows that the congruences~\eqref{eq:cong} hold unconditionally for all even $m\ge2$, $i>i_m$ and $n\ge0$. Moreover $\nu\big(\OPT_m(2^i)-\OPT_m(2^{i-1})\big)=3i+c_m$ exactly.
\end{corollary}

For example, $(c_2,c_4,c_8,c_{16},c_{32},c_{64})=(-4,-1,1,3,5,7)$. The first three values agree with the data of~\cite{ST}.

\subsection*{Idea of the proof}
We work in the variable $t=x-1$. Section~\ref{sec:closed} gives closed forms for the images $\Th(t^k)$, obtained from the substitution $x=(y+y^{-1})/2$. They show that $\Th(t^k)$ is divisible by $2^{k+1}$ and has very few terms of degree $<2$ in $t$. Consequently the pair $(R(1),R'(1))$ is transformed by an explicit linear recursion, modulo the right power of $2$. In this form the conjectured pattern becomes an invariant of the recursion (Proposition~\ref{prop:step}), and it is started by explicit base cases (Section~\ref{sec:base}). A Lean 4 formalization is described in Appendix~\ref{app:lean}.

\section{Preliminaries}\label{sec:prelim}

\subsection{Chebyshev polynomials}
Let $y$ be an indeterminate. Since $(y+y^{-1})/2$ is transcendental over $\Q$, the assignment $x\mapsto (y+y^{-1})/2$ defines an injective ring homomorphism $\iota\colon\Q[x]\to\Q(y)$. We use three standard facts (see e.g.\ \cite{MH}).
\begin{enumerate}[(T1)]
\item $\iota(T_n)=(y^n+y^{-n})/2$. This follows by induction from the recurrence, since $(y+y^{-1})(y^n+y^{-n})-(y^{n-1}+y^{1-n})=y^{n+1}+y^{-n-1}$.
\item $T_{2j}(x)=T_j(2x^2-1)$. Indeed $\iota(2x^2-1)=(y^2+y^{-2})/2$, and (T1) applied with $y^2$ in place of $y$ gives the claim, by injectivity of $\iota$.
\item $T_n(1)=1$, $T_n'(1)=n^2$, and $T_n''(1)=n^2(n^2-1)/3$. By (T1), $T_n(\cos\theta)=\cos n\theta$, so $(1-x^2)T_n''-xT_n'+n^2T_n=0$. Setting $x=1$ gives $T_n'(1)=n^2T_n(1)=n^2$. Differentiating once more and setting $x=1$ gives $-3T_n''(1)+(n^2-1)T_n'(1)=0$.
\end{enumerate}

\subsection{The operators \texorpdfstring{$\Sop$}{S} and \texorpdfstring{$\Th$}{Theta}}
Since $T_n(1)=1$, we have $\Sop(F)(1)=F(1)$ for every $F\in\Z[x]$. In particular $\Sop((x-1)R)$ vanishes at $x=1$, and as $x-1$ is monic, $\Th(R)\in\Z[x]$. By linearity and (T1),
\begin{equation}\label{eq:Siota}
  \iota\big(\Sop(F)\big)=\tfrac12\big(F(y^2)+F(y^{-2})\big)\qquad(F\in\Q[x]),
\end{equation}
where $F(y^{\pm2})$ means $F$ evaluated at $x=y^{\pm2}$. Also, $R_{i+1}=\Th(R_i)$ for $i\ge i_m$, because $P_{i+1}-P_i=\Sop(P_i-P_{i-1})=\Sop((x-1)R_i)$ (this is \cite[Lemma~6.1]{ST}).

\subsection{The variable \texorpdfstring{$t=x-1$}{t = x - 1}}
The substitution $F(x)\mapsto F(1+t)$ is a ring isomorphism $\Z[x]\to\Z[t]$ with inverse $G(t)\mapsto G(x-1)$. Both maps send $2^r\Z[\cdot]$ into $2^r\Z[\cdot]$. Hence, if we write
\[
  R(1+t)=\sum_{k\ge0}r_kt^k ,
\]
then $\nu(R)=\min_k\nu(r_k)$. In other words, $\nu$ may be computed in the $t$-basis. Moreover $r_0=R(1)$ and $r_1=R'(1)$. From now on we identify $F\in\Z[x]$ with $F(1+t)\in\Z[t]$, and write $\Th(t^k)$ for $\Th((x-1)^k)$ expressed in $t$. Under $\iota$,
\begin{equation}\label{eq:tiota}
  \iota(t)=\iota(x-1)=\frac{(y-1)^2}{2y},\qquad \iota(t+2)=\iota(x+1)=\frac{(y+1)^2}{2y}.
\end{equation}

\section{The action of \texorpdfstring{$\Th$}{Theta} on powers of \texorpdfstring{$t$}{t}}\label{sec:closed}

\begin{lemma}\label{lem:yform}
For every $k\ge0$,
\[
  \iota\big(\Th(t^k)\big)=y\,(y-1)^{k-1}(y+1)^{k+1}\big(1+(-1)^{k+1}y^{-2k-2}\big).
\]
\end{lemma}
\begin{proof}
Take $F(x)=(x-1)^{k+1}$ in~\eqref{eq:Siota}. Then $F(y^2)=(y^2-1)^{k+1}$ and $F(y^{-2})=(y^{-2}-1)^{k+1}=(-1)^{k+1}y^{-2k-2}(y^2-1)^{k+1}$, so
$\iota\big(\Sop((x-1)^{k+1})\big)=\tfrac12(y^2-1)^{k+1}\big(1+(-1)^{k+1}y^{-2k-2}\big)$.
Dividing by $\iota(x-1)=(y-1)^2/(2y)$ gives the claim.
\end{proof}

Put $\omega:=2t(t+2)\in\Z[t]$, and for $l\ge0$ let
\[
  W_l(u):=\frac{T_{l+1}(1+u)-T_l(1+u)}{u}\in\Z[u].
\]
This is a polynomial with integer coefficients, because $T_{l+1}(1+u)-T_l(1+u)\in\Z[u]$ vanishes at $u=0$ by (T3). For example $W_0=1$ and $W_1(u)=3+2u$.

\begin{lemma}[Closed forms]\label{lem:closed}
In $\Z[t]$ we have
\begin{align*}
  \Th(t^{2l-1})&=4^l\,t^{l-1}(t+2)^l\,T_l(1+\omega) &&(l\ge1),\\
  \Th(t^{2l})&=2\cdot4^l\,t^{l}(t+2)^{l+1}\,W_l(\omega) &&(l\ge0).
\end{align*}
\end{lemma}
\begin{proof}
By~\eqref{eq:tiota}, $\iota(\omega)=(y^2-1)^2/(2y^2)$ and $\iota(1+\omega)=(y^2+y^{-2})/2$. Hence (T1) gives $\iota(T_l(1+\omega))=(y^{2l}+y^{-2l})/2$. Since $\iota$ is injective, it suffices to compare images under~$\iota$ with Lemma~\ref{lem:yform}.

\emph{Odd exponents, $k=2l-1$.} The image of the right-hand side is
\begin{align*}
  4^l\,\frac{(y-1)^{2l-2}}{(2y)^{l-1}}\,\frac{(y+1)^{2l}}{(2y)^{l}}\,\frac{y^{2l}+y^{-2l}}{2}
  &=y^{1-2l}(y-1)^{2l-2}(y+1)^{2l}\big(y^{2l}+y^{-2l}\big)\\
  &=y(y-1)^{2l-2}(y+1)^{2l}\big(1+y^{-4l}\big).
\end{align*}
This is Lemma~\ref{lem:yform} with $k=2l-1$.

\emph{Even exponents, $k=2l$.} First,
\begin{align*}
  \iota\big(T_{l+1}(1+\omega)-T_l(1+\omega)\big)&=\tfrac12\big(y^{2l+2}+y^{-2l-2}-y^{2l}-y^{-2l}\big)\\
  &=\tfrac12(y^2-1)\big(y^{2l}-y^{-2l-2}\big),
\end{align*}
so $\iota(W_l(\omega))=y^2\big(y^{2l}-y^{-2l-2}\big)/(y^2-1)$. Hence the image of the right-hand side is
\[
  2\cdot4^l\,\frac{(y-1)^{2l}(y+1)^{2l+2}}{(2y)^{2l+1}}\cdot\frac{y^2\big(y^{2l}-y^{-2l-2}\big)}{(y-1)(y+1)}
  =y(y-1)^{2l-1}(y+1)^{2l+1}\big(1-y^{-4l-2}\big).
\]
This is Lemma~\ref{lem:yform} with $k=2l$.
\end{proof}

\begin{lemma}\label{lem:ab}
\begin{enumerate}[(a)]
\item $\Th(t^k)\in 2^{k+1}\Z[t]$ for every $k\ge0$.
\item We have $\Th(1)=4+2t$ and $\Th(t)=8+36t+32t^2+8t^3$. Moreover
\[
  \Th(t^2)\equiv 96t,\qquad \Th(t^3)\equiv 64t,\qquad \Th(t^k)\equiv0\quad(k\ge4)\pmod{t^2}.
\]
\end{enumerate}
\end{lemma}
\begin{proof}
(a) In Lemma~\ref{lem:closed} the factors $t$, $t+2$, $T_l(1+\omega)$ and $W_l(\omega)$ have integer coefficients. The prefactor is $4^l=2^{k+1}$ when $k=2l-1$, and $2\cdot 4^l=2^{k+1}$ when $k=2l$.

(b) From Lemma~\ref{lem:closed}:
\begin{itemize}
\item $\Th(1)=2(t+2)$, using $W_0=1$.
\item $\Th(t)=4(t+2)(1+4t+2t^2)=8+36t+32t^2+8t^3$.
\item $\Th(t^2)=8t(t+2)^2(3+2\omega)\equiv 8t\cdot4\cdot3=96t \pmod{t^2}$, using $W_1(u)=3+2u$ and $\omega\equiv0\pmod t$.
\item $\Th(t^3)=16\,t(t+2)^2\,T_2(1+\omega)\equiv 16t\cdot 4\cdot T_2(1)=64t \pmod{t^2}$.
\item For $k=2l-1\ge5$ the closed form contains the factor $t^{l-1}$ with $l-1\ge2$. For $k=2l\ge4$ it contains $t^l$ with $l\ge2$. Either way $\Th(t^k)\equiv0\pmod{t^2}$.\qedhere
\end{itemize}
\end{proof}

\begin{corollary}\label{cor:rec}
Let $R=\sum_k r_kt^k\in\Z[t]$ and write $\Th(R)=\sum_kr'_kt^k$. Then
\[
  r'_0=4r_0+8r_1,\qquad r'_1=2r_0+36r_1+96r_2+64r_3 .
\]
\end{corollary}
\begin{proof}
Apply Lemma~\ref{lem:ab}(b) to $\Th(R)=\sum_kr_k\Th(t^k)$. The first identity is the formula $\Th(R)(1)=4R(1)+8R'(1)$ of \cite[Lemma~6.1]{ST}.
\end{proof}

\section{An invariant}\label{sec:inv}

\begin{definition}\label{def:good}
Let $v\ge0$. A polynomial $R=\sum_kr_kt^k\in\Z[t]$ (equivalently, an element of $\Z[x]$) is \emph{good of level $v$} if
\begin{enumerate}[(G1)]
\item $\nu(R)=v$;
\item $\nu(r_1)=v$, that is, $\nu(R'(1))=v$;
\item $\nu(r_0)=v+1$, that is, $\nu(R(1))=v+1$;
\item $\alpha+\beta\equiv2\pmod 4$, where $\alpha:=r_0/2^{v+1}$ and $\beta:=r_1/2^{v}$.
\end{enumerate}
\end{definition}
By (G2) and (G3) the integers $\alpha$ and $\beta$ are odd. Condition (G4) therefore says that $\alpha\equiv\beta\pmod4$. Conditions (G1)--(G3) are exactly the three assertions of Conjecture~\ref{conj} at a single index. Condition (G4) is the extra information needed to propagate them: by Corollary~\ref{cor:rec}, the new constant term is $4r_0+8r_1=2^{v+3}(\alpha+\beta)$. Since $\alpha$ and $\beta$ are odd, its valuation is at least $v+4$, and (G4) makes it exactly $v+4$. This is the cancellation between $4R(1)$ and $8R'(1)$ observed in \cite[\S6]{ST}.

\begin{proposition}[Inductive step]\label{prop:step}
If $R$ is good of level $v$, then $\Th(R)$ is good of level $v+3$.
\end{proposition}
\begin{proof}
Put $\rho_k:=r_k/2^v$. These are integers by (G1), and $\rho_0=2\alpha$, $\rho_1=\beta$. By Corollary~\ref{cor:rec},
\begin{align*}
  r'_0&=4r_0+8r_1=2^{v+3}(\alpha+\beta),\\
  r'_1&=2r_0+36r_1+96r_2+64r_3=2^{v+2}\gamma,\qquad \gamma:=\alpha+9\beta+24\rho_2+16\rho_3 .
\end{align*}
Since $\alpha+\beta\equiv2\pmod4$, we get $\nu(r'_0)=v+4$. Since $\gamma\equiv\alpha+\beta\pmod 8$, we have $\nu(\gamma)=1$ and hence $\nu(r'_1)=v+3$.

Next we show that $\Th(R)\in 2^{v+3}\Z[t]$. Write $\Th(R)=r_0\Th(1)+r_1\Th(t)+\sum_{k\ge2}r_k\Th(t^k)$. For $k\ge2$, Lemma~\ref{lem:ab}(a) gives $r_k\Th(t^k)\in2^{v+k+1}\Z[t]\subseteq 2^{v+3}\Z[t]$. By Lemma~\ref{lem:ab}(b),
\[
  r_0\Th(1)+r_1\Th(t)=(4r_0+8r_1)+(2r_0+36r_1)\,t+32r_1\,t^2+8r_1\,t^3 .
\]
The coefficients here are $2^{v+3}(\alpha+\beta)$, $2^{v+2}(\alpha+9\beta)$, $2^{v+5}\beta$ and $2^{v+3}\beta$. They are all divisible by $2^{v+3}$, because $\alpha+9\beta$ is even. Hence $\nu(\Th(R))\ge v+3$, and equality holds because $\nu(r'_1)=v+3$. This proves (G1)--(G3) at level $v+3$.

Finally, $\alpha'=r'_0/2^{v+4}=(\alpha+\beta)/2$ and $\beta'=r'_1/2^{v+3}=\gamma/2$. Therefore
\[
  \alpha'+\beta'=\frac{2\alpha+10\beta+24\rho_2+16\rho_3}{2}=\alpha+5\beta+12\rho_2+8\rho_3\equiv\alpha+\beta\equiv2\pmod4,
\]
which is (G4) at level $v+3$.
\end{proof}

\begin{lemma}[Base case B1]\label{lem:B1}
If $R\in\Z[x]$ and $R(1)$ is odd, then $\Th(R)$ is good of level~$1$.
\end{lemma}
\begin{proof}
By Corollary~\ref{cor:rec}, $r'_0=4(r_0+2r_1)$ and $r'_1=2(r_0+18r_1+48r_2+32r_3)$. Since $r_0$ is odd, $\nu(r'_0)=2$ and $\nu(r'_1)=1$. By Lemma~\ref{lem:ab}(a), $\Th(R)\in2\Z[t]$, so $\nu(\Th(R))=1$. Finally,
\[
  \alpha'+\beta'=(r_0+2r_1)+(r_0+18r_1+48r_2+32r_3)\equiv2r_0\equiv2\pmod 4.\qedhere
\]
\end{proof}

\section{The base cases}\label{sec:base}

\begin{lemma}[$m\equiv2\pmod4$]\label{lem:base2}
Let $m\equiv2\pmod 4$ and $n=m/2$, which is odd. Then $R^{(m)}_2(1)=n^2(4n^2+5)/3$, which is odd. Hence $R^{(m)}_3$ is good of level~$1$.
\end{lemma}
\begin{proof}
Here $i_m=2$ and $R_2=(\Sop(T_n)-T_n)/(x-1)$. If $G\in\Z[x]$ and $G(1)=0$, then $\big(G/(x-1)\big)(1)=G'(1)$. For $F=\sum_jf_jx^j$ we have
\[
  \Sop(F)'(1)=\sum_jf_jT_{2j}'(1)=4\sum_jj^2f_j=4\big(F''(1)+F'(1)\big)
\]
by (T3). With $F=T_n$ and (T3), this gives
\[
  R_2(1)=4\Big(\frac{n^2(n^2-1)}3+n^2\Big)-n^2=\frac{n^2(4n^2+5)}{3}.
\]
The numerator is odd, so the integer $R_2(1)$ is odd. Lemma~\ref{lem:B1} applied to $R_3=\Th(R_2)$ finishes the proof.
\end{proof}

\begin{lemma}[$m\equiv4\pmod 8$]\label{lem:base4}
Let $m\equiv4\pmod8$ and $n=m/4$, which is odd. Then $R^{(m)}_1(1)=n(4n-1)$, which is odd. Hence $R^{(m)}_2$ is good of level~$1$.
\end{lemma}
\begin{proof}
Here $i_m=1$ and $R_1=(T_{2n}-x^n)/(x-1)$, so $R_1(1)=T_{2n}'(1)-n=4n^2-n$. Apply Lemma~\ref{lem:B1}.
\end{proof}

\begin{lemma}[$8\mid m$; base case B2]\label{lem:base8}
Let $8\mid m$. Put $n=m/4$, which is even, and $N=2n=m/2$. Then $R^{(m)}_2$ is good of level $v:=2\nu(n)+1=2\nu(m)-3$.
\end{lemma}
\begin{proof}
Put $Q:=(T_N-1)/(x-1)$ and $G:=(x^n-1)/(x-1)$, both in $\Z[x]$. Then $R_1=(T_N-x^n)/(x-1)=Q-G$. Moreover
\[
  \Th(G)=\frac{\Sop(x^n-1)}{x-1}=\frac{T_{2n}-1}{x-1}=Q,
\]
so $R_2=\Th(R_1)=\Th(Q)-Q$.

\emph{Step 1: $Q$ is good of level $v$.} Write $Q(1+t)=\sum_kq_kt^k$. Since $Q(1+t)=(T_N(1+t)-1)/t$, Taylor expansion and (T3) give
\[
  q_0=T_N'(1)=N^2=4n^2,\qquad q_1=\tfrac12T_N''(1)=\frac{N^2(N^2-1)}6=\frac{2n^2(4n^2-1)}3 .
\]
Thus $\nu(q_0)=2\nu(n)+2=v+1$ and $\nu(q_1)=2\nu(n)+1=v$.

To show $\nu(Q)\ge v$, put $F_M:=T_M(1+t)-1\in\Z[t]$ for $M\ge1$. The identity $T_{2M}=2T_M^2-1$ follows from (T1), since $2\big((y^M+y^{-M})/2\big)^2-1=(y^{2M}+y^{-2M})/2$. It gives
\[
  F_{2M}=2F_M\,(F_M+2).
\]
Since $\nu(FG)\ge\nu(F)+\nu(G)$ for $F,G\in\Z[t]$, the bound $\nu(F_M)\ge c$ implies $\nu(F_{2M})\ge 1+c+\min(c,1)$. We claim that $\nu(F_M)\ge2\nu(M)-1$ whenever $\nu(M)\ge1$, and argue by induction on $\nu(M)$. If $M=2M'$ with $M'$ odd, then $\nu(F_{M'})\ge0$, so $\nu(F_M)\ge1$. If $M=2M'$ with $\nu(M')\ge1$, then $c:=2\nu(M')-1\ge1$ and $\nu(F_M)\ge c+2=2\nu(M)-1$. Since $F_N=t\cdot Q(1+t)$, we obtain $\nu(Q)=\nu(F_N)\ge 2\nu(N)-1=v$. Together with $\nu(q_1)=v$, this gives $\nu(Q)=v$.

It remains to check (G4) for $Q$. Write $n=2^{\nu(n)}n'$ with $n'$ odd. Then $\alpha_Q=q_0/2^{v+1}=n'^2\equiv1\pmod 8$ and $\beta_Q=q_1/2^{v}=n'^2\cdot(4n^2-1)/3$. Since $n$ is even, $3\cdot\frac{4n^2-1}{3}\equiv-1\pmod{16}$, so $(4n^2-1)/3\equiv5\pmod{16}$. Hence $\beta_Q\equiv5\pmod8$, and $\alpha_Q+\beta_Q\equiv6\equiv2\pmod4$. So $Q$ is good of level $v$.

\emph{Step 2.} By Proposition~\ref{prop:step}, $\Th(Q)$ is good of level $v+3$. In particular $\Th(Q)\in2^{v+3}\Z[t]$, and therefore
\[
  R_2\equiv -Q\pmod{2^{v+3}\Z[t]} .
\]
Since $v+1<v+3$, it follows that $\nu(R_2)=\nu(Q)=v$, $\nu(r_1)=\nu(q_1)=v$ and $\nu(r_0)=\nu(q_0)=v+1$, where $R_2(1+t)=\sum r_kt^k$. Furthermore $\alpha_{R_2}\equiv-\alpha_Q\pmod4$ and $\beta_{R_2}\equiv-\beta_Q\pmod 8$, so $\alpha_{R_2}+\beta_{R_2}\equiv-2\equiv2\pmod4$. Hence $R_2$ is good of level~$v$.
\end{proof}

\section{Proofs of Theorem~\texorpdfstring{\ref{thm:main}}{1.2} and Corollary~\texorpdfstring{\ref{cor:cm}}{1.3}}\label{sec:proof}

The theorem concerns only the polynomials $R^{(m)}_i$. To pass to the congruences we need the $q$-series identity~\eqref{eq:65}, which is \cite[(6.5)]{ST}. To keep the note self-contained we include a short proof of it. Let $\varphi(q)=\sum_{k\in\Z}q^{k^2}$ and $(a;q)_\infty=\prod_{j\ge0}(1-aq^j)$.

\begin{lemma}[{cf.\ \cite[(2.16), (6.3)]{ST}}]\label{lem:Uiden}
\begin{enumerate}[(a)]
\item $A^2=\varphi(q)/\varphi(-q)$, and $\xi=A^4\in1+8q\Z[[q]]$.
\item $U(A^{2k})=T_k(\xi)$ for every $k\ge0$. Consequently $U\big(F(\xi)\big)=\Sop(F)(\xi)$ for every $F\in\Z[x]$.
\item For every even $m\ge2$ and $i\ge i_m-1$ we have $U^i(A^m)=P_i(\xi)$. Hence~\eqref{eq:65} holds for $i\ge i_m$.
\end{enumerate}
\end{lemma}
\begin{proof}
(a) By the Jacobi triple product identity \cite[Theorem~2.8]{And}, $\varphi(q)=(q^2;q^2)_\infty(-q;q^2)_\infty^2$, and replacing $q$ by $-q$ gives $\varphi(-q)=(q^2;q^2)_\infty(q;q^2)_\infty^2$. Since $A=(-q;q^2)_\infty/(q;q^2)_\infty$, this gives $A^2=\varphi(q)/\varphi(-q)$. Next, $\varphi(q)-\varphi(-q)=4\sum_{k\ge1,\ k\text{ odd}}q^{k^2}\in4q\Z[[q]]$, and $\varphi(-q)\in1+q\Z[[q]]$ is a unit of $\Z[[q]]$. Hence $A^2-1\in4q\Z[[q]]$, so $A^2+1\in2\Z[[q]]$ and $\xi-1=(A^2-1)(A^2+1)\in8q\Z[[q]]$.

(b) We use two classical identities:
\begin{equation}\label{eq:phi}
  \varphi(q)\varphi(-q)=\varphi(-q^2)^2,\qquad \varphi(q)^2+\varphi(-q)^2=2\varphi(q^2)^2 .
\end{equation}
For the first, note that $(-q;q^2)_\infty(q;q^2)_\infty=(q^2;q^4)_\infty$ and $(q^2;q^2)_\infty=(q^2;q^4)_\infty(q^4;q^4)_\infty$. By the product formulas in (a), both sides then equal $(q^4;q^4)_\infty^2(q^2;q^4)_\infty^4$. For the second, write $\varphi(q)^2=\sum_n r(n)q^n$, where $r(n)$ is the number of $(a,b)\in\Z^2$ with $a^2+b^2=n$. Then $\varphi(q)^2+\varphi(-q)^2=2\sum_n r(2n)q^{2n}$. If $a^2+b^2=2n$ then $a\equiv b\pmod 2$, and $(a,b)\mapsto\big(\tfrac{a+b}2,\tfrac{a-b}2\big)$ is a bijection onto the solutions of $c^2+d^2=n$, with inverse $(c,d)\mapsto(c+d,c-d)$. So $r(2n)=r(n)$.

Put $X=A^2$, viewed in the field $\Q((q))$. For $F\in\Z[[q]]$ we have $U(F)(q^2)=\tfrac12\big(F(q)+F(-q)\big)$. By (a), $X(-q)=X(q)^{-1}$. By (T1), the identity $T_k\big((z+z^{-1})/2\big)=(z^k+z^{-k})/2$ holds in $\Q[z,z^{-1}]$, and we may specialise it to $z=X$. Therefore
\[
  U(X^k)(q^2)=\frac{X^k+X^{-k}}2=T_k\Big(\frac{X+X^{-1}}2\Big),
\]
and by~\eqref{eq:phi}
\[
  \frac{X+X^{-1}}2=\frac{\varphi(q)^2+\varphi(-q)^2}{2\varphi(q)\varphi(-q)}=\frac{\varphi(q^2)^2}{\varphi(-q^2)^2}=X(q^2)^2=\xi(q^2).
\]
Hence $U(A^{2k})(q^2)=T_k(\xi(q^2))$, that is, $U(A^{2k})=T_k(\xi)$. Taking $k=2j$ gives $U(\xi^j)=T_{2j}(\xi)$, and the second assertion follows by linearity of $U$.

(c) If $4\mid m$, then $U^0(A^m)=\xi^{m/4}=P_0(\xi)$. If $m\equiv2\pmod4$, then (b) with $k=m/2$ gives $U(A^m)=T_{m/2}(\xi)=P_1(\xi)$. In both cases the start is $i=i_m-1$, and (b) gives $U(P_i(\xi))=\Sop(P_i)(\xi)=P_{i+1}(\xi)$. Finally, for $i\ge i_m$,
\[
  \Dc^{(m)}_i=U^i(A^m)-U^{i-1}(A^m)=P_i(\xi)-P_{i-1}(\xi)=(\xi-1)R^{(m)}_i(\xi).\qedhere
\]
\end{proof}

\begin{proof}[Proof of Theorem~\ref{thm:main}]
Let $i_0:=i_m+1$. By Lemmas~\ref{lem:base2}, \ref{lem:base4} and~\ref{lem:base8}, $R_{i_0}$ is good of some level $\lambda_m$:
\begin{itemize}
\item if $m\equiv2\pmod4$, then $i_0=3$ and $\lambda_m=1=3\cdot3-8$;
\item if $m\equiv4\pmod8$, then $i_0=2$ and $\lambda_m=1=3\cdot2+2\cdot2-9$;
\item if $8\mid m$, then $i_0=2$ and $\lambda_m=2\nu(m)-3=3\cdot 2+2\nu(m)-9$.
\end{itemize}
Since $R_{i+1}=\Th(R_i)$ for $i\ge i_m$, Proposition~\ref{prop:step} and induction show that $R_i$ is good of level $\lambda_m+3(i-i_0)$ for every $i\ge i_0$. This level equals the value of $v_i$ claimed in the theorem. For each $i>i_m$, condition (G1) gives $v_i=\nu(R_i)$, (G2) gives $\nu(R_i'(1))=v_i$, and (G3) gives $\nu(R_i(1))=v_i+1$. The equality $v_{i+1}=v_i+3$ is the shift of the level in Proposition~\ref{prop:step}.
\end{proof}

\begin{proof}[Proof of Corollary~\ref{cor:cm}]
We recall the argument of \cite[proof of Cor.~6.3]{ST}, which uses only~\eqref{eq:65}, $\xi\in1+8q\Z[[q]]$ (both proved in Lemma~\ref{lem:Uiden}) and the three assertions of Conjecture~\ref{conj}. Fix $i>i_m$ and write $R_i=2^{v_i}R^\ast_i$ with $R^\ast_i\in\Z[x]$. By Theorem~\ref{thm:main}, $R^\ast_i(1)=R_i(1)/2^{v_i}\equiv 2\pmod 4$. Since $R^\ast_i(x)-R^\ast_i(1)\in(x-1)\Z[x]$, we get $R^\ast_i(\xi)-R^\ast_i(1)\in(\xi-1)\Z[[q]]\subseteq8\Z[[q]]$. So $R^\ast_i(\xi)\in2\Z[[q]]$, and by~\eqref{eq:65} $\Dc_i^{(m)}=2^{v_i}(\xi-1)R^\ast_i(\xi)\in2^{v_i+4}\Z[[q]]$. Moreover, since $\xi-1=8q+O(q^2)$, the coefficient of $q$ in $\Dc^{(m)}_i$ equals $8R_i(1)$, which has valuation exactly $v_i+4$. Thus~\eqref{eq:cong} holds with $3i+c_m=v_i+4$, and the modulus is attained at $n=1$. Inserting the values of $v_i$ from Theorem~\ref{thm:main} gives $c_m=(3i-8)+4-3i=-4$ when $m\equiv2\pmod4$, and $c_m=(3i+2\nu(m)-9)+4-3i=2\nu(m)-5$ when $4\mid m$.
\end{proof}

\section{Remarks}\label{sec:remarks}

\begin{remark}[The hypothesis $i>i_m$]\label{rem:im}
The hypothesis $i>i_m$ in Conjecture~\ref{conj} cannot be dropped in general. Let $i=i_m$. Then $v_{i_m}=\nu(R_{i_m})=0$ in all cases. For $m\not\equiv0\pmod8$ this holds because $R_{i_m}(1)$ is odd by Lemmas~\ref{lem:base2} and~\ref{lem:base4}. For $8\mid m$ it holds because $R_1=Q-G$, where $\nu(Q)\ge3$ and $G=1+x+\dots+x^{n-1}$.
\begin{itemize}
\item If $m\not\equiv0\pmod 8$, then $R_{i_m}(1)$ is odd, so $\nu(R_{i_m}(1))=0\neq v_{i_m}+1$. For example, $R^{(2)}_2=2x+1$, and $R^{(12)}_1=(T_6-x^3)/(x-1)$ satisfies $R^{(12)}_1(1)=33$. Moreover $v_{i_m+1}=1\ne v_{i_m}+3$.
\item If $16\mid m$, then $R_1'(1)=q_1-n(n-1)/2$, and $\nu(q_1)=2\nu(n)+1>\nu(n)-1=\nu\big(n(n-1)/2\big)$. Hence $\nu(R_1'(1))=\nu(m)-3\ge1\ne v_1$. For example, $m=16$ gives $(\nu(R_1),\nu(R_1'(1)),\nu(R_1(1)))=(0,1,2)$.
\item If $m\equiv8\pmod{16}$, the same computation gives $\nu(R_1'(1))=0$ and $\nu(R_1(1))=\nu(4n^2-n)=1$, and $v_2=3$ by Lemma~\ref{lem:base8}. So all three assertions happen to hold at $i=i_m$ as well.
\end{itemize}
A computation for all even $m\le400$ confirms this trichotomy. The congruence~\eqref{eq:cong} itself can also fail at $i=i_m$. Since $A=1+2q+2q^2+O(q^3)$, we have $\OPT_m(1)=2m$ and $\OPT_m(2)=2m^2$. Hence $\OPT_m(2)-\OPT_m(1)=2m(m-1)$ has valuation $\nu(m)+1$. For $4\mid m$ this is smaller than $3i_m+c_m=2\nu(m)-2$ exactly when $16\mid m$. For instance $\nu(\OPT_{16}(2)-\OPT_{16}(1))=\nu(480)=5<6$.
\end{remark}

\begin{remark}
The proof gives slightly more than the conjecture asks for. For $i>i_m$, the normalised values $\alpha_i=R_i(1)/2^{v_i+1}$ and $\beta_i=R_i'(1)/2^{v_i}$ are odd and satisfy $\alpha_i\equiv\beta_i\pmod 4$. Moreover $\alpha_{i+1}=(\alpha_i+\beta_i)/2$. The case distinction in $c_m$ between $\nu(m)=1$ and $\nu(m)\ge2$ reflects the different starting index $i_m$. For $m\equiv2\pmod 4$ the first good polynomial $R_3$ has level~$1$. For $4\mid m$ the first good polynomial $R_2$ has level $2\nu(m)-3$, which equals $1$ when $\nu(m)=2$.
\end{remark}

\section{Computational verification}\label{sec:verif}

The proof above is complete and does not rely on computation. As a safeguard, all statements were also checked by computer. The checks used several programs, written separately and using exact integer arithmetic, or arithmetic modulo $2^K$ with $K$ far above every valuation involved. The software was PARI/GP~2.15.4~\cite{PARI} and Python~3 with python-flint~0.9.0~\cite{FLINT} and NumPy.
\begin{enumerate}[1.]
\item \emph{The conjecture and the formula for $v_i$, from the definitions of \cite[\S6]{ST}:}
\begin{itemize}
\item all even $m\le400$ with $i\in\{i_m+1,i_m+2\}$ (PARI/GP);
\item all even $m\le200$ with $i\le i_m+5$ (exact modulo $2^{160}$, with $\Sop$ evaluated by Clenshaw summation);
\item all even $m\le200$ with $i\le i_m+6$, and all even $m\le100$ with $i\le i_m+8$ (modulo $2^{64}$, an independent implementation);
\item all even $m\le200$ with $i\le i_m+3$ (exact big integers, python-flint; this check also confirms $R_{i+1}=\Th(R_i)$ exactly);
\item $m\in\{2^k:k\le11\}\cup\{3\cdot2^k:k\le9\}\cup\{5\cdot2^k:k\le8\}$ with $i\le i_m+3$ (modulo $2^{64}$), which covers $\nu(m)$ up to~$11$.
\end{itemize}
No discrepancy was found.
\item \emph{The lemmas:}
\begin{itemize}
\item the closed forms of Lemma~\ref{lem:closed} and Lemma~\ref{lem:ab}(b), for $k\le60$;
\item Lemma~\ref{lem:ab}(a), for $k\le80$; in fact the exact content of $\Th(t^k)$ is $2^{k+1}$;
\item the parity statements of Lemmas~\ref{lem:base2} and~\ref{lem:base4}, for $m\le2400$;
\item Lemma~\ref{lem:base8} (including the residues $\alpha_Q\equiv1$, $\beta_Q\equiv5\pmod 8$), for $8\mid m\le1200$.
\end{itemize}
\item \emph{The congruences themselves, directly from the generating function.} We used the identity $A^2=\varphi(q)/\varphi(-q)$, where $\varphi(q)=\sum_{k\in\Z}q^{k^2}$. We verified $\nu\big(\OPT_m(2^in)-\OPT_m(2^{i-1}n)\big)\ge3i+c_m$, with equality at $n=1$, in the following ranges:
\begin{itemize}
\item all even $m\le128$, $i_m<i\le14$ and $2^in\le 16384$ (800 pairs $(m,i)$);
\item separately, from the eta-quotient $f_2^3/(f_1^2f_4)$, for even $m\le16$, $i\le7$ and $2^in\le6000$.
\end{itemize}
\end{enumerate}
All scripts, together with the Lean source file of Appendix~\ref{app:lean}, are included as ancillary files with the arXiv version of this note (directory \texttt{anc/}).

\section*{AI usage statement}
The research presented here is the result of human--AI interaction involving Claude (Anthropic). The author selected the problem, directed the work, and decided what to check and how to present it. The proof of Theorem~\ref{thm:main}, including the change of variable $t=x-1$, the closed forms of Section~\ref{sec:closed}, the invariant of Section~\ref{sec:inv} and the base cases, was found by the AI system in the course of this interaction. The same holds for the proof of Lemma~\ref{lem:Uiden}, the verification programs of Section~\ref{sec:verif} and the Lean formalization of Appendix~\ref{app:lean}. The manuscript was drafted with the AI system and revised under the author's direction.

The results were checked in three independent ways: by separate AI review passes that re-derived each step, by the computations of Section~\ref{sec:verif} (programs written independently of one another, see the ancillary files), and by the Lean~4 formalization of Theorem~\ref{thm:main}, which uses no unproved assumptions. The note has not yet been reviewed by a human specialist in the area. Conjecture~\ref{conj} was posed in~\cite{ST} as open, supported by computations for $m\le14$, and to the author's knowledge it was unsolved before this note. The author takes full responsibility for the content and welcomes detailed questions about the role of AI in this work.

\appendix
\section{A Lean 4 formalization}\label{app:lean}

Theorem~\ref{thm:main} has been formalized and machine-checked in Lean~4 (toolchain \texttt{v4.35.0-rc4}) with Mathlib (the matching tag \texttt{v4.35.0-rc4}), in a single file \texttt{SaikiaTalukdar.lean}, which is included as the ancillary file \texttt{anc/lean/SaikiaTalukdar.lean}. The definitions follow \cite[\S6]{ST} literally:
\begin{itemize}
\item the Chebyshev polynomials $D_n\in\Z[x]$, by the three-term recurrence;
\item the $\Z$-linear operator $\Sop$ with $\Sop(x^j)=D_{2j}$;
\item $i_m$ and $P_{i_m-1}$ as in the introduction, with $P_i=\Sop^{\,i-i_m+1}(P_{i_m-1})$;
\item $R_i=(P_i-P_{i-1})/(x-1)$, computed as division by the monic polynomial $x-1$;
\item $\nu(F)$, the minimum of the $2$-adic valuations of the nonzero coefficients of $F$.
\end{itemize}
The file proves the following.
\begin{itemize}
\item \texttt{conjecture\_6\_2}: for every even $m\ge2$ and every $i>i_m$, $R_i\ne0$, $R_i(1)\neq0$ and $R_i'(1)\ne0$; moreover $\nu(R_i'(1))=\nu(R_i)$, $\nu(R_i(1))=\nu(R_i)+1$ and $\nu(R_{i+1})=\nu(R_i)+3$.
\item \texttt{nu\_R\_of\_mod4\_eq2} and \texttt{nu\_R\_of\_four\_dvd}: the two formulas for $v_i$ in Theorem~\ref{thm:main}.
\end{itemize}
The formal proof follows Sections~\ref{sec:closed}--\ref{sec:proof}: the $t$-coordinate, the bounds of Lemma~\ref{lem:ab}, the invariant of Section~\ref{sec:inv}, and the three base cases, including the content argument based on $F_{2M}=2F_M(F_M+2)$. The only difference is that the closed forms of Lemma~\ref{lem:closed} are obtained by a forward-difference argument rather than the substitution $x=(y+y^{-1})/2$. For each of the three theorems, the command \texttt{\#print axioms} reports only \texttt{propext}, \texttt{Classical.choice} and \texttt{Quot.sound}. In particular, no \texttt{sorry} is used. Corollary~\ref{cor:cm} depends on the $q$-series identities of \cite{ST}, and it is not formalized.

\begin{thebibliography}{FLINT}
\bibitem[And]{And} G.~E. Andrews, \emph{The Theory of Partitions}, Encyclopedia of Mathematics and its Applications~2, Addison--Wesley, Reading, MA, 1976; reprinted by Cambridge University Press, 1998.
\bibitem[FLINT]{FLINT} The FLINT team, \emph{FLINT: Fast Library for Number Theory}, and \emph{python-flint}, \url{https://flintlib.org}.
\bibitem[MH]{MH} J.~C. Mason and D.~C. Handscomb, \emph{Chebyshev Polynomials}, Chapman \& Hall/CRC, Boca Raton, 2003.
\bibitem[PARI]{PARI} The PARI Group, \emph{PARI/GP version 2.15.4}, Univ.\ Bordeaux, \url{https://pari.math.u-bordeaux.fr}.
\bibitem[ST]{ST} M.~P. Saikia and P.~Talukdar, \emph{Internal congruences modulo powers of $2$ for overpartition tuples with odd parts}, arXiv:2609.11806v1 [math.NT], 10~Sep 2026.
\end{thebibliography}

\end{document}
